\metatitle{Hivert polynomials}
\metadescription{An introduction to Hivert polynomials, their definition as quasisymmetric analogs of Hall--Littlewood polynomials, and their properties including expansions and connections to other quasisymmetric functions.}
\metakeywords{Hivert polynomials,quasisymmetric functions,Hall--Littlewood polynomials,Gessel quasisymmetric functions,monomial quasisymmetric functions}

\section[hivert]{Hivert polynomials}

\begin{polydata}{hivert}
  Name   & Hivert polynomials \\
  Space    & QSym \\
  Basis    & True \\
  Rating   & 1 \\
  Bib      & Hivert2000 \\
  Year     & 2000 \\
  Symbol   & $\hivert_\alpha(\xvec;t)$ \\
  Category & Schur \\
\end{polydata}


Hivert polynomials were introduced by \name{Florent Hivert} in \cite{Hivert2000},
as a quasisymmetric analog of the \hyperref[hallLittlewoodP]{Hall--Littlewood polynomials}.
These interpolate between the \hyperref[gessel]{Gessel quasisymmetric functions}
and the \hyperref[qmonom]{monomial quasisymmetric functions}.

Transition matrices involving Hivert polynomials are studied in \cite{LoehrSerranoWarrington2013}.

Let $\alpha=(\alpha_1,\dotsc,\alpha_k)$ be a composition and let $\beta$ be a
refinement of $\alpha$.  If the parts of $\beta$ refining $\alpha_j$ are
counted by $b_j$, set
\[
\mathrm{Bre}(\beta,\alpha)=(b_1,\dotsc,b_k), \qquad
s(\alpha,\beta)\coloneqq \sum_{j=1}^k j(b_j-1).
\]
One convenient normalization of the \defin{Hivert polynomial} is
\[
\hivert_\alpha(\xvec;t)
=
\sum_{\beta\text{ refines }\alpha}
(-1)^{\ell(\beta)-\ell(\alpha)}
t^{s(\alpha,\beta)}\gessel_\beta(\xvec).
\]
This is \cite[Thm.~11]{LoehrSerranoWarrington2013}, written as a definition.
It gives
\[
\hivert_\alpha(\xvec;0)=\gessel_\alpha(\xvec),
\qquad
\hivert_\alpha(\xvec;1)=\qmonom_\alpha(\xvec).
\]

\begin{example}
% Related Rust: sym-poly/qsym/src/hivert.rs and
% sym-poly/qsym/examples/hivert_site_example.rs
For $\alpha=(3)$, the refinements are $(3)$, $(2,1)$, $(1,2)$, and
$(1,1,1)$.  Thus
\[
\hivert_3
=
\gessel_3
-t\gessel_{21}
-t\gessel_{12}
+t^2\gessel_{111}.
\]
In the monomial quasisymmetric basis this becomes
\[
\hivert_3
=
\qmonom_3
+(1-t)\qmonom_{21}
+(1-t)\qmonom_{12}
+(1-t)^2\qmonom_{111}.
\]
The inverse transition begins
\[
\gessel_3
=
\hivert_3
+t\hivert_{21}
+t\hivert_{12}
+t^3\hivert_{111},
\]
which is the degree-three case of \cite[Thm.~26]{LoehrSerranoWarrington2013}.
\end{example}

Hivert also introduced quasisymmetric divided-difference operators.  Replacing
the ordinary divided differences by Hivert's operators in the construction of
Schur functions recovers the \hyperref[gessel]{fundamental quasisymmetric
functions}.  \name{Angela Hicks} and \name{Elizabeth Niese}
\cite{HicksNiese2024x} apply the same idea to key polynomials and Demazure
atoms.  The resulting analogues are the \hyperref[slideF]{fundamental slide
polynomials} and the \hyperref[fundamentalParticle]{fundamental particle
basis}, respectively.
